\paragraph{The normalization of alpha and its signed configuration.}
\label{inc-ampl-clausura}
The coordinate of $\alpha$ is determined along the positional route by
composing the regional blocks and the dodecaphasic register. Before
positional division, the oriented combination has components
$Z_j=P_j+E_j-\Phi_j-K_j^{\mathrm{per}}$, and the balance with the incoming quotient
is $s_j=Z_j+c_{j+1}$. The normalization satisfies
\begin{equation}
 A_j=s_j-1000c_j=Z_j+c_{j+1}-1000c_j,\qquad 0\leq A_j<1000.
 \label{inc-ampl-normalizacion}
\end{equation}
This identity preserves both the normalized block and the
contribution of the transported quotient. The limit of the compatible
prefixes determines the scalar coordinate of the closure. The signed
configuration, in turn, retains which positions of the combination
before normalization exhibit a negative defect.
In the initial window, this configuration is the vector $Z_0$ in
\eqref{exc:precarry}, and its negative support is $H^-_\alpha$.

This defines two kinds of content of a single operation. Positional
normalization determines a real coordinate through its compatible
remainders and quotients; the negative support determines a
configuration of positions in the twelve-component cycle. The sign
refers to the register balance before normalization. Its incidence
readout remains available because the state retains that balance
together with the normalized result. This articulation belongs to the
proved positional route; the comparison with the analytic vacancy
operators retains the precise scope established in the chapter on
$\alpha$.
